\clearpage
\fancyhead[L]{\small INSTRUMENTE.}
\hypertarget{native-427-part-lower}{}
\begin{center}
\Large INSTRUMENTE.
\end{center}



\begin{center}
Göttingische gelehrte Anzeigen. Stück 167. Seite 1665..1669. 18. Oktober 1819.
\end{center}





Im 127. Stück dieser Anzeigen vom vorigen Jahr [S. 410 dieses Bandes] haben wir den Meridiankreis von \textsc{Repsold} beschrieben und einige der ersten Beobachtungen mitgeteilt, die damit an der hiesigen Sternwarte angestellt wurden. Jetzt berichten wir kurz über das \emph{Mittagsfernrohr} von \textsc{Reichenbach}, das zweite Hauptinstrument, dessen Aufstellung vor Kurzem abgeschlossen wurde. Da die Einrichtung der Mittagsfernrohre, wie sie beim gegenwärtigen Stand der Beobachtungskunst verwendet werden, in allen wesentlichen Punkten stets dieselbe ist, können wir uns hier auf die Besonderheiten unseres Instruments beschränken. Das Fernrohr hat eine Brennweite von sechs Pariser Fuß und eine Öffnung von 52 Linien. In dieser Hinsicht wird es also nur vom neuen Mittagsfernrohr der Sternwarte in Greenwich übertroffen, das bei 10 Fuß Brennweite eine Öffnung von 5 englischen Zoll hat. Die Achse misst 37 Zoll, und die Stahlzapfen ruhen in Pfannen aus Glockenmetall. Es gibt 4 Okulare, die nach einer vorläufigen Bestimmung 80-, 110-, 150- und 210-fach vergrößern. Auch die stärkste Vergrößerung liefert vollkommen scharfe Bilder. Herr Hofrat \textsc{Gauss} verwendet sie gewöhnlich, obwohl sie eine etwas sorgfältigere Einstellung der Höhe verlangt, weil das Gesichtsfeld dabei nur 9 Minuten im Durchmesser hat. Das Fadennetz besteht aus Spinnenfäden, von denen sieben senkrecht stehen. Sterne am Äquator durchlaufen die Zwischenräume in 10 Sekunden. Hat man die Fertigkeit, die das Beobachten bei diesen kurzen Intervallen verlangt, einmal erworben, so ist es überaus angenehm, in kurzer Zeit so viele Fadendurchgänge beobachten zu können. Bei der größeren Zahl von Fäden fällt es außerdem weniger ins Gewicht, wenn ein Durchgang an einem oder einigen Fäden absichtlich oder zufällig unbeobachtet bleibt. Aus mehreren hundert eigenen Beobachtungen von Sternen, die nicht zu weit vom Äquator entfernt sind, findet Herr Hofrat \textsc{Gauss} für den sogenannten wahrscheinlichen Fehler des beobachteten Durchgangs an einem Faden $=0^{\mathrm s}095$. An der Libelle, mit der die Achse waagerecht gestellt wird, entspricht ein Ausschlag von einem Pariser Zoll einer Neigung von 20 Sekunden. Das Fernrohr wird auf jede vorgegebene Deklination mit einem kleinen Kreis eingestellt, der am Fernrohr selbst nahe dem Okularende befestigt und unmittelbar in Viertelgrade sowie durch den Nonius in Minuten unterteilt ist. Auf ihm dreht sich eine Alhidade mit einer kleinen Libelle, und der Index an der Alhidade zeigt sofort die Deklination an. Diese Einrichtung hat den großen Vorteil, dass man das Instrument auch in umgekehrter Lage mit gleicher Leichtigkeit auf jede Höhe richten kann. Ein zweiter ähnlicher Kreis, den der Instrumentenbauer noch nachliefern wird und der auf der anderen Seite anzubringen ist, bietet dann den weiteren Vorteil, zwei Sterne, die schnell nacheinander durch die Meridianebene gehen, ohne Zeitverlust beobachten zu können. ---
Die optische Vollkommenheit des Fernrohrs ist bewundernswert. Einige Beispiele, die wir hier aus den Beobachtungen von Herrn Hofrat \textsc{Gauss} anführen, mögen zeigen, was man mit guten, aber nicht gerade ungewöhn%
\SourceLeaf{428}%
lich scharfen Augen in unserem Klima bei günstigem, aber keineswegs außergewöhnlichem Zustand der Luft mit gebührender Aufmerksamkeit, doch ohne übermäßige Anstrengung erkennen kann. Bei \emph{voller Fadenbeleuchtung} ist ein Saturnstrabant, häufig sind auch zwei, gut zu erkennen und zu beobachten. Dasselbe gilt für den Begleitstern des Polarsterns. Der kleine Stern, der jetzt nur 5 Minuten vom Pol entfernt ist, verträgt noch genügend Beleuchtung, um sich gut beobachten zu lassen. Schon in der Dämmerung, wenn die Fäden noch gar keine Beleuchtung brauchen, lässt er sich erkennen und beobachten. Bei Tage lassen sich Sterne bis zur 3. Größe meist leicht beobachten, wenn sie weder zu tief noch zu nahe bei der Sonne stehen. Bei günstiger Luft und unter Beachtung einiger Vorsichtsregeln sind auch Sterne vierter, ja selbst fünfter Größe noch bei Tage recht gut zu erkennen. So hat Herr Hofrat \textsc{Gauss} zum Beispiel die Sterne fünfter Größe $\gamma'$ im Kleinen Bären, $e$ im Großen Bären, $1$ im Drachen nach Hevel und $g$ im Großen Bären (Alcor) bei hellem Tage beobachtet, den letzten sogar zur Mittagszeit. Die Schwierigkeit, derart lichtschwache Sterne bei Tage zu beobachten, besteht eigentlich nur darin, dass die Bauart des Instruments kein ganz genaues Einstellen zulässt. Wo \emph{diese} Schwierigkeit wegfällt und das Auge also nur einen sehr kleinen Teil des Gesichtsfelds fixieren muss, erkennt man Sterne 5ter Größe, ja selbst noch lichtschwächere, leicht. Das ist bei Doppelsternen der Fall. So ist etwa der Begleiter von $\zeta$ im Großen Bären, dem die Sternverzeichnisse nur die sechste Größe zuweisen, am hellen Mittag sehr schön zu sehen. Ebenso ist der noch lichtschwächere Begleiter von $\alpha$ im Herkules gegenwärtig, da er nachmittags um 4 Uhr kulminiert, sehr gut zu erkennen. Im Übrigen liegen die vielen Vorteile der Möglichkeit, auch lichtschwächere Sterne bei Tage gut zu beobachten, für Kenner auf der Hand. Wir erwähnen hier nur einen sehr wichtigen. Bei der Berichtigung des Instruments anhand von Zirkumpolarsternen, die bekanntlich als einzige einem solchen Instrument angemessen ist und so oft wie möglich vorgenommen werden muss, bleibt man nicht auf den Polarstern beschränkt, dessen eine Kulmination oft in eine unbequeme Nachtstunde fällt oder dessen Beobachtung durch ungünstige Umstände vereitelt wird.

\par\indent Die Vorrichtung zum Ausheben und Umlegen des Instruments verdient noch besondere Erwähnung. Das Fernrohr wird mit einem Seil aus den Lagern gehoben, das oben über eine Rolle läuft und bis zu einer unten an der Mauer befestigten Winde geführt ist. Dann wird das Fernrohr frei in der Luft schwebend ohne jede Erschütterung umgewandt und wieder eingelegt. Der ganze Vorgang dauert nur wenige Minuten. Die über dem Fernrohr der Decke des Zimmers befindliche Rolle wird nach getaner Arbeit durch einen Zug zur Seite geschoben, damit sie beim Beobachten von Zenitalsternen nicht im Weg ist. Diesen ganzen schönen Apparat hat der hiesige Herr Inspektor \textsc{Rumpf} sehr zweckmäßig und solide ausgeführt.

\par\indent Als Beispiele fügen wir noch einige der ersten mit diesem Instrument angestellten Beobachtungen bei. Wir bemerken nur, dass die Deklinationen diesmal nicht beobachtet werden konnten, weil der Meridiankreis von \textsc{Repsold} in diesem Monat abgenommen war.

\Needspace{27\baselineskip}
\begin{center}
\emph{Beobachtungen des Saturn.}

\small
\setlength{\tabcolsep}{5pt}
\begin{tabular}{l|rrr|rrr}
\multicolumn{1}{c|}{1819}  &  \multicolumn{3}{c|}{Mittlere Zeit}  &  \multicolumn{3}{c}{Scheinbare Rektaszension}  \\
\hline
Sept.\ 4  &  $13^{\mathrm u}$  &  $5^{\mathrm m}$  &  $43^{\mathrm s}9$  &  $359^\circ$  &  $40'$  &  $2''7$  \\
5  &  $13$  &  $1$  &  $22.0$  &   &  $36$  &  $2.7$  \\
6  &  $12$  &  $57$  &  $10.7$  &   &  $32$  &  $9.6$  \\
7  &  $12$  &  $52$  &  $58.8$  &   &  $28$  &  $10.3$  \\
8  &  $12$  &  $48$  &  $46.1$  &   &  $24$  &  $4.0$  \\
9  &  $12$  &  $44$  &  $34.4$  &   &  $20$  &  $0.4$  \\
10  &  $12$  &  $40$  &  $22.2$  &   &  $15$  &  $55.6$  \\
11  &  $12$  &  $36$  &  $10.1$  &   &  $11$  &  $51.7$  \\
12  &  $12$  &  $31$  &  $57.5$  &   &  $7$  &  $41.2$  \\
\hline
 \multicolumn{1}{c|}{1819}  &  \multicolumn{3}{c|}{Mittlere Zeit}  &  \multicolumn{3}{c}{Scheinbare Rektaszension}  \\
\hline
 Sept.\ 13  &  $12^{\mathrm u}$  &  $27^{\mathrm m}$  &  $45^{\mathrm s}0$  &  $359^\circ$  &  $3'$  &  $30''4$  \\
 14  &  $12$  &  $23$  &  $32.3$  &  $358$  &  $59$  &  $18.1$  \\
 15  &  $12$  &  $19$  &  $19.8$  &   &  $55$  &  $9.0$  \\
 16  &  $12$  &  $15$  &  $7.0$  &   &  $50$  &  $54.4$  \\
 19  &  $12$  &  $2$  &  $28.6$  &   &  $38$  &  $10.6$  \\
 21  &  $11$  &  $54$  &  $3.0$  &   &  $29$  &  $42.9$  \\
 24  &  $11$  &  $41$  &  $24.2$  &   &  $15$  &  $56.4$  \\
 28  &  $11$  &  $24$  &  $33.4$  &   &  $0$  &  $5.1$  \\
 Okt.\ 1  &  $11$  &  $11$  &  $55.9$  &  $357$  &  $47$  &  $36.3$  \\
\end{tabular}
\end{center}

\Needspace{22\baselineskip}
\SourceLeaf{429}
\emph{Beobachtungen der Vesta.}

\begin{center}
\begin{tabular}{l|rrr|rrr}
\multicolumn{1}{c|}{1819}  &  \multicolumn{3}{c|}{Mittlere Zeit}  &  \multicolumn{3}{c}{Scheinbare Rektaszension}  \\
\hline
Sept.\ 7  &  $13^{\mathrm u}$  &  $32^{\mathrm m}$  &  $52^{\mathrm s}0$  &  $9^\circ$  &  $28'$  &  $7''2$  \\
9  &   &  $23$  &  $32.6$  &   &  $6$  &  $10.5$  \\
10  &   &  $18$  &  $51.1$  &  $8$  &  $54$  &  $44.7$  \\
11  &   &  $14$  &  $8.6$  &   &  $43$  &  $1.8$  \\
12  &   &  $9$  &  $25.1$  &   &  $31$  &  $6.9$  \\
15  &  $12$  &  $55$  &  $18.7$  &  $7$  &  $53$  &  $49.8$  \\
\hline
 \multicolumn{1}{c|}{1819}  &  \multicolumn{3}{c|}{Mittlere Zeit}  &  \multicolumn{3}{c}{Scheinbare Rektaszension}  \\
\hline
 Sept.\ 16  &  $12^{\mathrm u}$  &  $50^{\mathrm m}$  &  $21^{\mathrm s}4$  &  $7^\circ$  &  $40'$  &  $58''5$  \\
 19  &   &  $35$  &  $55.5$  &   &  $1$  &  $18.0$  \\
 21  &   &  $26$  &  $15.9$  &  $6$  &  $34$  &  $15.9$  \\
 24  &   &  $11$  &  $43.4$  &  $5$  &  $52$  &  $58.0$  \\
 28  &  $11$  &  $52$  &  $18.1$  &  $4$  &  $57$  &  $24.3$  \\
 Okt.\ 1  &   &  $37$  &  $45.3$  &   &  $16$  &  $1.2$  \\
\end{tabular}
\end{center}

\Needspace{27\baselineskip}
\par\indent Da diese Beobachtungen noch nicht mit Tafeln oder Elementen verglichen sind, geben wir auch noch die Beobachtungen eines Fixsterns samt der Reduktion auf den mittleren Ort für den Anfang des Jahres 1820 an.


\begin{center}
\emph{Rektaszension von 85 Pegasi.}
\end{center}

\begin{center}
\begin{tabular}{l|rrr|rrr}
\multicolumn{1}{c|}{1819}  &  \multicolumn{3}{c|}{Beobachtete}  &  \multicolumn{3}{c}{Mittlere für 1820}  \\
\hline
Sept.\ 4  &  $358^\circ$  &  $11'$  &  $40''8$  &  $358^\circ$  &  $11'$  &  $43''0$  \\
5  &   &   &  $41.4$  &   &   &  $43.4$  \\
6  &   &   &  $40.2$  &   &   &  $42.0$  \\
7  &   &   &  $43.9$  &   &   &  $45.5$  \\
8  &   &   &  $42.6$  &   &   &  $44.0$  \\
9  &   &   &  $42.9$  &   &   &  $44.1$  \\
10  &   &   &  $43.0$  &   &   &  $44.1$  \\
\hline
 \multicolumn{1}{c|}{1819}  &  \multicolumn{3}{c|}{Beobachtete}  &  \multicolumn{3}{c}{Mittlere für 1820}  \\
\hline
 Sept.\ 11  &  $358^\circ$  &  $11'$  &  $42''6$  &  $358^\circ$  &  $11'$  &  $43''6$  \\
 12  &   &   &  $45.4$  &   &   &  $46.2$  \\
 13  &   &   &  $41.8$  &   &   &  $42.4$  \\
 14  &   &   &  $44.7$  &   &   &  $45.1$  \\
 15  &   &   &  $44.8$  &   &   &  $45.1$  \\
 16  &   &   &  $43.0$  &   &   &  $43.2$  \\
\multicolumn{4}{r|}{Im Mittel . .} &  $358$  &  $11$  &  $43.97$  \\
\end{tabular}
\end{center}

\par\indent Dieser Stern zeichnet sich durch eine starke Eigenbewegung aus. Der Vergleich der aus diesen Beobachtungen hervorgehenden Rektaszension mit der von \textsc{Bradley} für 1755 ergibt für 64.7 Jahre eine Eigenbewegung von $+54''65$, also für ein Jahr $+0''845$.
